% !PRINTED 381
\weberpage{381}
%% [running head, not body text -- commented out 2026-08-07] Page furniture (running
%% head, folio, signature) is excluded from this edition by standing policy.
%% §\,104. Charaktere und Geschlechter der quadratischen Formen. \hfill 381

Haben $\delta_1, \delta_2$ keinen gemeinschaftlichen Teiler, so ist diese
symbolische Multiplikation eine wirkliche. Haben sie aber einen
gemeinschaftlichen Teiler, so ist noch ein quadratischer Faktor
aus dem wahren Produkt $\delta_1\delta_2$ abzuwerfen. Setzen wir in (2)
$\delta_1 = \delta_2$, so ist $\delta = 1$, d. h. es ist jedes Element der Gruppe
der $\delta$ sich selbst reziprok, und es besteht zugleich mit (2):

\weq{3}{\delta_1 = \delta\delta_2.}

Die Elemente dieser Gruppe lassen sich folgendermaßen durch eine
\gesperrt{Basis} darstellen (Bd. II, § 11):

\weq{4}{\delta = \delta_1^{\varepsilon_1}\delta_2^{\varepsilon_2}\dots\delta_\lambda^{\varepsilon_\lambda},}

worin $\varepsilon_1, \varepsilon_2, \dots \varepsilon_\lambda$ die Werte 0 oder 1 haben, und $\lambda$ die oben
angegebene Bedeutung hat. $\delta_1, \delta_2, \dots \delta_\lambda$ sind in dem Falle 1) die
in $D$ aufgehenden ungeraden Primdiskriminanten $\pm p$, im Falle 2)
kommt dazu noch $-4$, im Falle 3) $8$, im Falle 4) $-8$, im
Falle 5) $8$ und $-8$ (oder $-4$ und $8$ oder $-4$ und $-8$).

Mit Hilfe der Stammteiler werden nun die \gesperrt{Charaktere} der
primitiven quadratischen Formen definiert.

Wir wählen in einer Formenklasse $A$ einen Repräsentanten
$(a,b,c)$, in dem $a$ positiv und relativ prim zu $D$ ist. Ist dann eine
zu $D$ teilerfremde Zahl $m$ durch $A$ darstellbar, so ist
\[
m = ax^2 + bxy + cy^2,
\]
\weq{5}{4am = (2ax + by)^2 - Dy^2.}

Es sei nun $\delta$ ein Stammteiler von $D$ und $m$ \gesperrt{relativ prim}
zu $\delta$, so ist, zunächst für ein ungerades $\delta$, nach § 85, (14), (16):

\weq{6}{(\delta, m) = (\delta, a).}

Diese Relation gilt aber auch für $\delta = -4\delta_1, +8\delta_1, -8\delta_1$.
Denn es ist in diesen Fällen $a$ und $m$ ungerade und nach (1)
und (5):
\[
\begin{aligned}
\delta &= -4\delta_1, & D &\equiv 0, -4 \ (\mathrm{mod}\ 16), & am &\equiv 1 \ (\mathrm{mod}\ 4), \\
\delta &= 8\delta_1, & D &\equiv 0, 8 \ (\mathrm{mod}\ 32), & am &\equiv \pm 1 \ (\mathrm{mod}\ 8), \\
\delta &= -8\delta_1, & D &\equiv 0, -8 \ (\mathrm{mod}\ 32), & am &\equiv 1, 3 \ (\mathrm{mod}\ 8),
\end{aligned}
\]
woraus nach § 85 auch für diese Fälle die Gleichung (6) folgt,
die danach allgemein bewiesen ist.

4. \gesperrt{Der Wert des Symbols $(\delta, m)$ ist also nicht von
der Zahl $m$, sondern nur von der Formenklasse $A$ abhängig,
durch die $m$ darstellbar ist. Es wird der zum Stammteiler $\delta$ gehörige Charakter dieser Klasse genannt und
mit $\chi(\delta, A)$ bezeichnet.}
